\documentclass[11pt]{article}

\usepackage[margin=1in]{geometry}
\usepackage{amsmath,amssymb,amsthm}
\usepackage{mathtools}
\usepackage{algorithm}
\usepackage{algpseudocode}
\usepackage{booktabs}
\usepackage{tabularx}
\usepackage{array}
\usepackage{xcolor}
\usepackage{microtype}
\usepackage[htt]{hyphenat}
\usepackage{hyperref}
\usepackage{enumitem}
\usepackage{listings}
\usepackage{tikz}
\usepackage{parskip}
\usepackage{titlesec}

% Prevent overfull hboxes from long \texttt{} identifiers and URLs.
\sloppy
\emergencystretch=3em
\hbadness=10000
\hfuzz=2pt

\hypersetup{
    colorlinks=true,
    linkcolor=blue!60!black,
    citecolor=blue!60!black,
    urlcolor=blue!60!black,
}

\titleformat{\section}
  {\normalfont\Large\bfseries}{\thesection}{1em}{}
\titleformat{\subsection}
  {\normalfont\large\bfseries}{\thesubsection}{1em}{}

\newcommand{\pmask}{p_{\text{mask}}}
\newcommand{\Ind}{\mathbb{1}}
\newcommand{\E}{\mathbb{E}}
\newcommand{\R}{\mathbb{R}}
\DeclareMathOperator{\Bern}{Bernoulli}
\DeclareMathOperator{\Cat}{Categorical}
\DeclareMathOperator{\Geom}{Geometric}
\DeclareMathOperator{\TopK}{TopK}

\title{\textbf{Masking Strategies for Antibody MLM Pretraining} \\
\large Algorithmic Design, Derivations, and Justifications}
\author{Ayan Goel}
\date{\today}

\begin{document}
\maketitle

\begin{abstract}
This document provides a comprehensive algorithmic description of every
masking strategy implemented in the antibody MLM pretraining framework.
Each strategy is presented with (i) its biological/statistical motivation,
(ii) the formal mathematical formulation, (iii) the concrete algorithm,
and (iv) justification of every design decision, including alternatives
that were considered and rejected. The document is deliberately focused
on algorithms, not on data, training, or evaluation.
\end{abstract}

\tableofcontents
\vspace{1em}

% ============================================================================
\section{Preliminaries and Common Framework}
% ============================================================================

\subsection{Notation}
Let $x = (x_1, x_2, \dots, x_L)$ denote a tokenized antibody sequence of
length $L$, with each $x_i$ an integer token ID from a protein vocabulary
$\mathcal{V}$. Let $s \in \{0, 1\}^L$ denote the \emph{special-tokens mask},
where $s_i = 1$ iff position $i$ is a special token
($[\texttt{CLS}]$, $[\texttt{SEP}]$, $[\texttt{PAD}]$, $[\texttt{MASK}]$, etc.).
Let $\mathcal{M}_{\text{eligible}} = \{i : s_i = 0\}$ be the set of
non-special positions that are eligible for masking, and let
$n = |\mathcal{M}_{\text{eligible}}|$.

A masking strategy defines a (possibly stochastic) map
\[
\mathcal{S} : (x, s, m) \;\mapsto\; (\tilde x, y),
\]
where $m$ is an optional metadata dictionary, $\tilde x$ is the corrupted
input, and $y \in (\mathcal{V} \cup \{-100\})^L$ are MLM targets with
$y_i = x_i$ at masked positions and $y_i = -100$ elsewhere (the standard
ignore-index for HuggingFace cross-entropy).

\subsection{Two-Stage Decomposition}
Every strategy in the framework implements exactly one function,
\texttt{select\_mask\_positions()}, returning a binary mask
$M \in \{0,1\}^L$. The base class handles a common second stage shared
across all strategies:

\paragraph{Stage 1: Position Selection.} Strategy-specific; chooses
$M_i = 1$ to indicate position $i$ is selected for masking.

\paragraph{Stage 2: Token Corruption.} For every position where $M_i = 1$,
the base class draws a Bernoulli-style corruption:
\begin{equation}
\tilde{x}_i \sim \begin{cases}
[\texttt{MASK}] & \text{w.p. } p_{\text{mask}} = 0.8, \\
\text{Uniform}(\mathcal{A}_{20}) & \text{w.p. } p_{\text{rand}} = 0.1, \\
x_i & \text{w.p. } p_{\text{keep}} = 0.1,
\end{cases}
\label{eq:bert-80-10-10}
\end{equation}
where $\mathcal{A}_{20} \subset \mathcal{V}$ is the set of token IDs
for the 20 canonical amino acids \texttt{ACDEFGHIKLMNPQRSTVWY}.

\paragraph{Why 80/10/10?} This is the BERT convention. It is preserved
because (i) it prevents the model from learning the trivial rule
``if the input is $[\texttt{MASK}]$, predict anything; otherwise copy,''
(ii) the random-replacement 10\% acts as a mild denoising objective,
and (iii) the 10\% keep ensures the model cannot ignore uncorrupted
positions.

\paragraph{Why restrict random replacements to canonical amino acids?}
The original HuggingFace reference implementation samples
\texttt{randint(low=5, high=vocab\_size)}, which can pick
non-standard tokens such as \texttt{X}, \texttt{B}, or \texttt{Z},
effectively replacing approximately $0.065\%$ of tokens with
biologically meaningless ``garbage'' characters. We precompute
$\mathcal{A}_{20}$ at strategy construction time:
\begin{equation}
\mathcal{A}_{20} \;=\; \{\text{tokenizer.convert\_tokens\_to\_ids}(a) : a \in \{\texttt{A},\texttt{C},\dots,\texttt{Y}\}\}
\end{equation}
and sample uniformly from it. This removes the garbage-token noise
source while still providing a stochastic denoising signal.

\paragraph{Stage 2 algorithm.} Given the selected mask $M$ and a sample
$u \sim \text{Uniform}[0,1]^{\sum_i M_i}$:
\begin{align}
\text{replace\_with\_mask}_j &\coloneqq \Ind[u_j < 0.8], \\
\text{replace\_with\_random}_j &\coloneqq \Ind[0.8 \le u_j < 0.9], \\
\text{keep}_j &\coloneqq \Ind[u_j \ge 0.9].
\end{align}
These define a partition of the $\sum_i M_i$ selected positions. The
keep-case intentionally leaves $\tilde x_i = x_i$; only the labels $y_i$
carry a gradient signal at those positions.

\subsection{The Central Equation: Budget-Normalized Weighted Bernoulli}
\label{sec:weighted-bernoulli}
All region- and label-aware strategies (\texttt{cdr}, \texttt{germline},
\texttt{interface}, parts of \texttt{multispecific}) share a single
position-selection primitive: a \emph{budget-normalized weighted Bernoulli}.
Given a per-position weight $w_i \ge 0$ with $w_i = 0$ at special tokens,
we define a per-position masking probability
\begin{equation}
\boxed{\;
p_i \;=\; \min\!\left(\, \pmask \cdot \frac{w_i}{\bar w}, \; 1 \right),
\quad\text{where}\quad \bar w = \frac{1}{n} \sum_{j \in \mathcal{M}_{\text{eligible}}} w_j,
\;}
\label{eq:weighted-bernoulli}
\end{equation}
and draw $M_i \sim \Bern(p_i)$ independently across positions.

\paragraph{Expected mask rate is preserved.} The defining property of
equation~(\ref{eq:weighted-bernoulli}) is that \emph{before} clamping,
the expected fraction of eligible positions masked equals $\pmask$:
\begin{equation}
\E\!\left[\frac{1}{n} \sum_{i \in \mathcal{M}_{\text{eligible}}} M_i\right]
\;=\; \frac{1}{n}\sum_{i} \pmask \cdot \frac{w_i}{\bar w}
\;=\; \pmask \cdot \frac{1}{\bar w} \cdot \frac{1}{n}\sum_i w_i
\;=\; \pmask.
\label{eq:budget-conservation}
\end{equation}
This is critical: it decouples the \emph{shape} of the mask distribution
(controlled by the weights $w_i$) from its \emph{total budget}
(controlled by $\pmask$). Without this normalization, increasing any
weight would drift the global masking rate up, confounding ablation
comparisons across strategies.

\paragraph{Why clamp?} For positions with $w_i / \bar w > 1/\pmask$
(e.g.\ extreme outliers), the raw scaled probability would exceed $1$.
Clamping prevents illegal probabilities but does slightly bias the
effective budget downward by an amount
\[
\Delta = \frac{1}{n}\sum_i \max\!\left(\pmask \cdot \tfrac{w_i}{\bar w} - 1,\; 0\right).
\]
In practice, with the chosen weights (max ratio $\approx 6{:}1$) and
$\pmask = 0.15$, clamping is essentially never triggered ($\Delta \approx 0$).

\paragraph{Why Bernoulli instead of multinomial top-$k$?} A multinomial
top-$k$ with $k = \lfloor \pmask \cdot n \rfloor$ would guarantee an exact
budget per sample but loses the natural variance of independent selection
and is more difficult to vectorize on GPU. The expected-budget Bernoulli
approach matches BERT-style uniform masking exactly when $w_i \equiv 1$,
so all strategies degenerate cleanly to the baseline.

% ============================================================================
\section{Strategy 1: Uniform Masking (Baseline)}
% ============================================================================

\subsection{Motivation}
This is the standard BERT masking policy and serves as the control
experiment against which all other strategies are compared. It assumes
nothing about the biological structure of the sequence --- every
non-special position is a priori equally informative.

\subsection{Formal Definition}
\begin{equation}
p_i = \pmask \cdot \Ind[s_i = 0], \qquad
M_i \sim \Bern(p_i).
\end{equation}

\subsection{Algorithm}
\begin{algorithm}[H]
\caption{Uniform Masking}
\begin{algorithmic}[1]
\State $P \gets \pmask \cdot \mathbf{1}_L$
\State $P[s = 1] \gets 0$
\State $M \sim \Bern(P)$
\State \Return $M$
\end{algorithmic}
\end{algorithm}

\subsection{Design Justification}
\begin{itemize}
\item \textbf{Why include at all?} Without a matched-budget uniform
baseline, observed gains from any biased strategy could be attributed
to (a) the biological prior, (b) changes in effective loss weighting, or
(c) merely a different random seed. The uniform baseline isolates
(a) from (b) and (c).
\item \textbf{Why independent Bernoulli instead of shuffling the first
$\lfloor \pmask n \rfloor$ positions?} Independent Bernoulli matches the
original BERT paper and the HuggingFace \texttt{DataCollatorForLanguageModeling}
behavior, preserving comparability with the broader MLM literature.
\item \textbf{Why $\pmask = 0.15$?} This is the BERT default. It is
preserved across all strategies so that any performance delta is not
confounded by a changed budget. Higher rates ($0.25$, $0.40$) were
considered because protein sequences have stronger local context, but
those experiments are kept as a separate ablation axis, not baked into
the strategies.
\end{itemize}

% ============================================================================
\section{Strategy 2: CDR-Focused Masking}
% ============================================================================

\subsection{Biological Motivation}
In antibody variable domains, the three Complementarity-Determining
Regions (CDR1, CDR2, CDR3) form the hypervariable loops that determine
antigen binding. Despite occupying only $\sim 25\%$ of residues, CDRs
carry the overwhelming majority of the sequence's functional information:
germline-derived framework residues are nearly constant, while CDR
residues vary enormously across antibodies. Uniform masking wastes
roughly $75\%$ of its learning budget on framework positions that any
decent model can already predict from context.

\subsection{Formal Definition}
Let $c \in \{0,1,2,3\}^L$ be the per-position region label with
$c_i = 0$ for framework, $c_i = k$ for CDR $k$. Define a weight vector
\begin{equation}
w^{(\text{cdr})}_i =
\begin{cases}
w_{\text{fr}} & c_i = 0 \text{ (framework)}, \\
w_{\text{c1}} & c_i = 1, \\
w_{\text{c2}} & c_i = 2, \\
w_{\text{c3}} & c_i = 3, \\
0 & s_i = 1 \text{ (special token)}.
\end{cases}
\end{equation}
Defaults are $(w_{\text{fr}}, w_{\text{c1}}, w_{\text{c2}}, w_{\text{c3}})
= (1, 3, 3, 6)$. Masking probabilities are computed via the
budget-normalized weighted Bernoulli of
Eq.~(\ref{eq:weighted-bernoulli}).

\subsection{Expected Mass Allocation}
Let $n_k = |\{i : c_i = k\}|$ and let $f_k = n_k / n$. Then the expected
fraction of the total mask budget landing in region $k$ is
\begin{equation}
\E\!\left[\frac{\#\{i : M_i = 1, c_i = k\}}{\#\{i : M_i = 1\}}\right]
\approx \frac{w_k f_k}{\sum_{k'} w_{k'} f_{k'}}.
\end{equation}
With typical $f = (0.75, 0.05, 0.07, 0.13)$ for VH and default weights,
this gives the CDR3 alone absorbing
\begin{equation}
\frac{6 \cdot 0.13}{1 \cdot 0.75 + 3 \cdot 0.05 + 3 \cdot 0.07 + 6 \cdot 0.13}
\approx \frac{0.78}{1.89} \approx 41\%,
\end{equation}
and CDR1+CDR2+CDR3 together absorbing roughly $60{-}70\%$ of masks
while being only $\sim 25\%$ of residues.

\subsection{Algorithm}
\begin{algorithm}[H]
\caption{CDR-Focused Masking}
\begin{algorithmic}[1]
\Require \texttt{cdr\_mask} $c \in \{0,1,2,3\}^L$
\If{$c$ is \textsc{None}} \Return \Call{UniformMasking}{$x, s$} \EndIf
\State $w_i \gets w^{(\text{cdr})}_{c_i}$ \Comment{vectorized lookup}
\State $w_i \gets 0$ for all $i$ with $s_i = 1$
\State $\bar w \gets \text{mean}\{w_i : s_i = 0\}$
\If{$\bar w > 0$}
  \State $P_i \gets \min\!\left(\pmask \cdot \tfrac{w_i}{\bar w},\, 1\right)$
\Else
  \State $P_i \gets \pmask$
\EndIf
\State $P_i \gets 0$ for all $i$ with $s_i = 1$
\State $M \sim \Bern(P)$
\State \Return $M$
\end{algorithmic}
\end{algorithm}

\subsection{Design Justification}
\begin{itemize}
\item \textbf{Why unequal CDR weights?} The three CDRs are not equally
informative. CDR3 is the only CDR with V(D)J junctional diversity and
carries the bulk of specificity determination; CDR1/CDR2 are
germline-encoded and less variable. Setting $w_{\text{c3}} = 2 w_{\text{c1,2}}$
concentrates gradient on the most informative positions without
starving the shorter loops entirely.
\item \textbf{Why keep $w_{\text{fr}} = 1$ instead of 0?} Setting
$w_{\text{fr}} = 0$ would prevent the model from ever predicting
framework residues, which are critical for VH--VL pairing and
structural stability. Keeping some framework mass trains the model to
recover the conserved Cys/Trp residues and prevents representation
collapse in framework tokens. The $1{:}3{:}3{:}6$ ratio was chosen
empirically as the simplest integer schedule that concentrates
$\approx 60\%$ mass in CDRs while retaining framework coverage.
\item \textbf{Why fall back to uniform when $c$ is unavailable?} CDR
annotation can fail (no OAS fields, ANARCI failure). A silent fallback
to uniform ensures training is never blocked by metadata issues and
degrades gracefully. The alternative --- raising an error or skipping
the sample --- biases the training set.
\item \textbf{Why the Bernoulli scheme instead of a hard rule
``always mask 2 CDR3 residues''?} A hard rule loses the
budget-conservation property and cannot compose with the hybrid
strategy's mixture sampler (where budgets from different sub-strategies
must be comparable).
\item \textbf{Rejected alternative: learned region weights.} Making
$w_k$ trainable parameters conditioned on $x$ couples the masker to the
model and complicates reproducibility. Fixed weights preserve a clean
ablation axis.
\end{itemize}

% ============================================================================
\section{Strategy 3: Span Masking (SpanBERT-style)}
% ============================================================================

\subsection{Motivation}
Standard token-level masking trivially leaks neighbor information:
predicting $x_i$ is easy when $x_{i-1}$ and $x_{i+1}$ are both visible,
because protein local context is strong (backbone hydrogen bonding,
$\alpha$-helix periodicity, $\beta$-strand register). SpanBERT showed
that masking contiguous spans forces the model to learn longer-range
dependencies and yields stronger representations.

In antibodies specifically, masking contiguous spans mimics the
real-world scenario of designing a complete sub-loop (e.g.\ an entire
CDR3 segment) rather than point mutations.

\subsection{Formal Definition}
Let $T = \max(1, \lfloor \pmask \cdot n \rceil)$ be the per-sequence
mask budget (in positions). The algorithm iteratively samples spans
until $T$ positions are covered. At each iteration:
\begin{enumerate}
\item Sample a start position $t$ uniformly from the currently-available
(non-special, non-masked) positions.
\item Sample a span length $\ell \sim \Geom(p_g)$, truncated to
$[1, \ell_{\max}]$. With $p_g = 0.2$, the mean span length is
\begin{equation}
\E[\ell] \;=\; \frac{1}{p_g} \;=\; 5 \quad (\text{truncated: } \approx 4.8).
\end{equation}
\item Mark positions $t, t+1, \dots, t+\ell-1$ as masked, \emph{stopping
early} at any special token or sequence boundary.
\end{enumerate}
The process stops when $\sum_i M_i \ge T$ or after $10T$ failed attempts.

\subsection{Geometric Span-Length Distribution}
The probability of sampling span length $\ell$ (before truncation) is
\begin{equation}
P(L = \ell) = (1 - p_g)^{\ell - 1}\, p_g, \qquad \ell = 1, 2, 3, \ldots
\end{equation}
With $p_g = 0.2$ and $\ell_{\max} = 10$ the truncation tail mass is
\begin{equation}
P(L > 10) = (1 - 0.2)^{10} \approx 0.107,
\end{equation}
i.e.\ roughly $11\%$ of raw samples get clipped to length $10$, so the
truncation is a soft cap rather than a hard constraint.

\subsection{Algorithm}
\begin{algorithm}[H]
\caption{Span Masking}
\begin{algorithmic}[1]
\State $T \gets \max(1, \text{round}(\pmask \cdot n))$
\State $M \gets \mathbf{0}_L$, $t_{\text{count}} \gets 0$, $\text{attempts} \gets 0$
\While{$t_{\text{count}} < T$ \textbf{and} $\text{attempts} < 10 T$}
  \State $\text{attempts} \gets \text{attempts} + 1$
  \State $A \gets \{i : s_i = 0 \wedge M_i = 0\}$
  \If{$A = \emptyset$} \textbf{break} \EndIf
  \State $t \sim \text{Uniform}(A)$
  \State $\ell \gets \min(\Geom(p_g) + 1,\, \ell_{\max})$
  \For{$o = 0$ to $\ell - 1$}
    \If{$t + o \ge L$ \textbf{or} $s_{t+o} = 1$} \textbf{break} \EndIf
    \If{$M_{t+o} = 0$}
      \State $M_{t+o} \gets 1$, $t_{\text{count}} \gets t_{\text{count}} + 1$
      \If{$t_{\text{count}} \ge T$} \textbf{break} \EndIf
    \EndIf
  \EndFor
\EndWhile
\State \Return $M$
\end{algorithmic}
\end{algorithm}

\subsection{Design Justification}
\begin{itemize}
\item \textbf{Why geometric distribution?} SpanBERT uses a truncated
geometric because (i) it has a clean single-parameter form,
(ii) it puts the most mass on short spans (capturing local motifs) but
still generates occasional long spans, and (iii) it induces a
tractable mean-length controllable by $p_g$.
\item \textbf{Why $p_g = 0.2$?} This is the SpanBERT default, giving
mean span length $\approx 5$. For a VH of length $\sim 120$, this means
roughly $T/5 = 0.15 \cdot 120 / 5 \approx 3.6$ spans per sequence ---
enough to sample multiple regions of the sequence per step, avoiding
a single-span regime that would be trivially localized.
\item \textbf{Why $\ell_{\max} = 10$?} An untruncated geometric
occasionally samples very long spans ($>$20) that would swallow entire
CDRs and leave the model with no local context at all. Capping at 10
(approximately one CDR3 length) preserves some within-region context
while still occasionally masking whole sub-loops.
\item \textbf{Why stop at special tokens?} Paired VH+VL sequences are
joined with a separator; a span that crosses the separator would mix
chains and corrupt the chain-boundary signal the model needs to learn
VH--VL pairing. Stopping-at-special is a hard invariant, not a
softened penalty.
\item \textbf{Why seed on available positions only?} Seeding uniformly
across \emph{all} positions (ignoring whether they are already masked)
would waste budget --- an already-masked seed re-attempts the same
neighborhood. Filtering on $(\neg s) \wedge (\neg M)$ keeps each attempt
productive.
\item \textbf{Why the $10T$ attempt cap?} A pathological failure mode
would be an infinite loop when $L$ is very small and almost everything
is already masked. The cap bounds the worst case.
\item \textbf{Rejected alternative: random-end span sampling.} Sampling
both ends independently and masking the interval $[\min, \max]$ was
considered but discarded: it produces a length distribution with
unclear statistics and breaks the controllability provided by
$\Geom(p_g)$.
\end{itemize}

% ============================================================================
\section{Strategy 4: Structure-Aware Masking}
% ============================================================================

\subsection{Motivation}
Masking a random subset of positions ignores 3D proximity. Two residues
that are 20 positions apart in sequence can be direct neighbors in the
folded structure (e.g.\ anti-parallel $\beta$-strands), and their joint
identity is correlated more strongly than sequence distance would
suggest. Structure-aware masking simulates the downstream task of
``fill in a surface patch'' by masking residues that cluster in 3D.

\subsection{Metadata Source}
Per-sequence spatial neighborhoods are precomputed offline via ESM2
attention contacts (see \texttt{scripts/predict\_structures.py}):
\begin{enumerate}
\item Run ESM2 on the raw sequence; extract the $L \times L$ contact
probability matrix from the symmetrized attention heads.
\item Zero out the diagonal (self is not a neighbor).
\item For each row $i$, take $\TopK$ with $K = 32$ (largest contact
probabilities).
\item Store as \texttt{knn\_indices} $\in \mathbb{Z}^{L \times K}$.
\end{enumerate}
This is strictly an offline step; training code never imports ESM2.

\subsection{Seed-and-Grow Algorithm}
Given a kNN graph, the strategy performs a simple seed-and-grow loop
to cover the budget with spatially-contiguous patches.

\begin{algorithm}[H]
\caption{Structure Masking: Seed-and-Grow}
\begin{algorithmic}[1]
\Require \texttt{knn\_indices} $N \in \mathbb{Z}^{L \times K}$
\State $\mathcal{M}_{\text{loc}} \gets$ non-special positions with non-empty kNN rows
\State $n' \gets |\mathcal{M}_{\text{loc}}|$; $T \gets \max(1, \text{round}(\pmask \cdot n'))$
\State Build a local index map: global $\to$ dense $[0, n')$
\State $M_{\text{loc}} \gets \mathbf{0}_{n'}$; $U_{\text{seed}} \gets \mathbf{0}_{n'}$; count $\gets 0$
\While{count $< T$}
  \State $S \gets \{j : U_{\text{seed},j} = 0 \wedge M_{\text{loc},j} = 0\}$
  \If{$S = \emptyset$} $S \gets \{j : M_{\text{loc},j} = 0\}$ \EndIf
  \If{$S = \emptyset$} \textbf{break} \EndIf
  \State $j^{\star} \sim \text{Uniform}(S)$
  \State $U_{\text{seed},j^{\star}} \gets 1$
  \If{$M_{\text{loc},j^{\star}} = 0$}
    \State $M_{\text{loc},j^{\star}} \gets 1$; count $\gets$ count $+ 1$
  \EndIf
  \For{$v$ in $N[j^{\star}, :]$ (closest first)}
    \If{count $\ge T$} \textbf{break} \EndIf
    \If{$M_{\text{loc},v} = 0$}
      \State $M_{\text{loc},v} \gets 1$; count $\gets$ count $+ 1$
    \EndIf
  \EndFor
\EndWhile
\State Map $M_{\text{loc}}$ back to global indices
\State \Return $M$
\end{algorithmic}
\end{algorithm}

\subsection{Design Justification}
\begin{itemize}
\item \textbf{Why $K = 32$?} ProteinMPNN, the state-of-the-art inverse
folding model, found that structural context saturates around $32$--$48$
nearest $C_\alpha$ neighbors; below that, the model lacks sufficient
context; above, the neighborhood includes distant residues whose
coupling is weak. $K = 32$ is a known sweet spot.
\item \textbf{Why seed-and-grow instead of single global cluster?}
Masking a \emph{single} 3D cluster per sequence would produce a
degenerate patch covering one face of the protein and leave the rest
of the fold intact. Multiple seeds sampled over training steps yield
broader coverage. Within a sequence, seed-and-grow also amortizes the
budget across several disjoint patches, avoiding the degenerate case
of a single cluster swallowing the entire CDR3.
\item \textbf{Why add closest-first neighbors instead of BFS?} Ordering
by contact probability ensures the most tightly-coupled residues are
masked first, maximizing the difficulty of the prediction task. A
plain BFS would add neighbors in order of graph discovery, which is
not physically meaningful.
\item \textbf{Why ESM2 contacts instead of AlphaFold/IgFold coordinates?}
ESM2 contacts are \emph{cheap} (forward pass, a few seconds per
sequence) while IgFold/AlphaFold require full structure prediction
(minutes per sequence) and do not scale to 500k sequences in a
reasonable wall-clock. The two are empirically correlated at the
contact-map level; the strategy degrades gracefully if one wants to
swap in a true Ca-coordinate source (\texttt{coords\_ca}) --- there
is a separate code path \texttt{\_mask\_from\_coords} that falls back
to Euclidean $C_\alpha$ distances and \texttt{torch.cdist}, kept for
backwards compatibility with the legacy IgFold sidecar.
\item \textbf{Why two code paths (\texttt{knn\_indices} vs
\texttt{coords\_ca})?} The \texttt{coords\_ca} path is a legacy IgFold
pipeline kept because small-scale experiments already have that
sidecar on disk. The \texttt{knn\_indices} path is the production
pipeline. If both are present, \texttt{knn\_indices} takes priority
because it is precomputed and avoids an $O(L^2)$ \texttt{cdist} at
every training step.
\item \textbf{Why fall back to uniform on missing metadata?} Training
must not crash because a single sequence lacks a structure
prediction. The fallback counter tracks how often this happens so
that a silent systematic failure (e.g.\ a broken preprocessing step)
surfaces in the logs.
\item \textbf{Why mark seeds as ``used'' but still let them be masked?}
A seed that is already in $M_{\text{loc}}$ does not re-contribute to
the count, but marking it in $U_{\text{seed}}$ ensures we don't loop
forever re-picking the same start. When no un-seeded positions remain,
the algorithm relaxes the constraint and lets any unmasked position
be a seed, ensuring progress.
\end{itemize}

% ============================================================================
\section{Strategy 5: Interface / Paratope Masking}
% ============================================================================

\subsection{Motivation}
The paratope is the subset of residues that make contact with the
antigen. These residues are, by definition, the functional output of
the antibody --- they are the positions whose identity determines
binding affinity and specificity. A masker that biases toward
paratope residues focuses learning on the decision surface of the
antibody.

\subsection{Metadata Source}
Paratope labels $\pi_i \in [0, 1]$ come from an empirical
$P(\text{paratope} \mid \text{AA}, \text{region})$ frequency table
(see \texttt{scripts/compute\_paratope\_labels.py}):
\begin{enumerate}
\item Take the TDC SAbDab\_Liberis dataset of antibody sequences
labeled with paratope positions (derived from crystal structures).
\item Annotate each TDC sequence with IMGT CDR regions via ANARCI.
\item Count $P(\text{paratope} \mid a, r)$ for each
(amino acid $a$, region $r \in \{\text{FR, CDR1, CDR2, CDR3}\}$) pair
with a minimum-count threshold of 5 (otherwise fall back to region-level
rate, then global rate).
\item For every training sequence, look up its per-residue
$\hat{\pi}_i = P(\text{paratope} \mid x_i, r_i)$.
\end{enumerate}
These \emph{soft} labels are floats in $[0,1]$, not hard 0/1, which is
handled natively by the weighted-Bernoulli formula.

\subsection{Formal Definition}
Given soft paratope labels $\pi \in [0,1]^L$ and weights
$(w_{\text{p}}, w_{\bar{\text{p}}}) = (6, 1)$:
\begin{equation}
w^{(\text{int})}_i = w_{\bar{\text{p}}} + (w_{\text{p}} - w_{\bar{\text{p}}}) \cdot \pi_i
= 1 + 5 \pi_i, \qquad \text{for } s_i = 0.
\end{equation}
Masking probabilities are then computed via
Eq.~(\ref{eq:weighted-bernoulli}).

\subsection{Linear Interpolation Justification}
The formula $w_i = w_{\bar{\text{p}}} + (w_{\text{p}} - w_{\bar{\text{p}}}) \pi_i$
is the unique linear map $[0,1] \to [w_{\bar{\text{p}}}, w_{\text{p}}]$
that matches the two endpoints exactly ($\pi_i = 0 \Rightarrow w_i = 1$,
$\pi_i = 1 \Rightarrow w_i = 6$) and behaves smoothly for intermediate
confidences. Any monotone map works; linear is the simplest and
introduces no hyperparameters beyond the endpoints.

\subsection{Algorithm}
\begin{algorithm}[H]
\caption{Interface Masking}
\begin{algorithmic}[1]
\Require \texttt{paratope\_labels} $\pi \in [0,1]^L$
\If{$\pi$ is \textsc{None}} \Return \Call{UniformMasking}{$x, s$} \EndIf
\State $w_i \gets w_{\bar{\text{p}}} + (w_{\text{p}} - w_{\bar{\text{p}}}) \pi_i$
\State $w_i \gets 0$ for $s_i = 1$
\State $\bar w \gets \text{mean}\{w_i : s_i = 0\}$
\If{$\bar w > 0$}
  \State $P_i \gets \min(\pmask w_i / \bar w, 1)$
\Else
  \State $P_i \gets \pmask$
\EndIf
\State $P_i \gets 0$ for $s_i = 1$
\State $M \sim \Bern(P)$; \Return $M$
\end{algorithmic}
\end{algorithm}

\subsection{Design Justification}
\begin{itemize}
\item \textbf{Why soft labels over hard 0/1?} Hard paratope calls
require a crystal structure; soft labels from
$P(\text{paratope} \mid a, r)$ can be computed for \emph{every}
training sequence without experimental structures. Soft labels also
let us weight positions by their \emph{expected} paratope probability
rather than a binary thresholded call.
\item \textbf{Why $(w_{\text{p}}, w_{\bar{\text{p}}}) = (6, 1)$?}
A ratio of $6{:}1$ concentrates approximately $50$--$60\%$ of the
mask budget onto paratope residues (which are only $\sim 15$--$25\%$
of the sequence) while leaving meaningful coverage on non-paratope
residues (so the model can still learn framework chemistry). The
same magnitude as the CDR3 weight keeps hyperparameter choices
consistent across strategies.
\item \textbf{Why compute the frequency table from TDC rather than
hardcode prior biochemical knowledge?} Hardcoded priors (e.g.\ ``Tyr,
Trp, Asn are paratope-prone'') are coarse and not region-specific.
The TDC-derived table is a per-(AA, region) estimate with empirical
support, no manual curation.
\item \textbf{Why the $\min$-count threshold of 5?} Rare $(a, r)$
combinations would have extreme point estimates dominated by noise.
Falling back to the region-level rate is a standard empirical-Bayes
shrinkage trick.
\item \textbf{Why fall back to uniform on missing metadata?} Same
argument as CDR: training robustness. Additionally, the strategy
tracks the fallback rate (\texttt{\_fallback\_calls / \_total\_calls})
and logs it periodically, so a broken sidecar file is detected
early.
\end{itemize}

% ============================================================================
\section{Strategy 6: Germline / SHM-Aware Masking}
% ============================================================================

\subsection{Motivation: The AbLang-2 Problem}
In adult human VH sequences, approximately $85\%$ of residues match
the germline V gene exactly; only $\sim 15\%$ have been mutated away
by somatic hypermutation (SHM) during affinity maturation. Uniform
masking therefore spends $85\%$ of its budget teaching the model to
predict \emph{germline} residues, which are easy (the V gene identity
largely determines them) and biologically uninteresting. The
functionally interesting residues are the SHM positions, which are
rare and get little gradient signal. AbLang-2 identified this as the
central bottleneck in antibody MLM pretraining.

\subsection{Metadata Source: Consensus-Based Germline Labels}
Per-residue germline labels $g_i \in \{0, 0.5, 1\}$ (and intermediate
values) are computed offline
(\texttt{scripts/compute\_germline\_labels.py}):
\begin{enumerate}
\item \textbf{Build germline reference.} For each V gene call
$v \in \{\text{IGHV1-2*01}, \ldots\}$, collect all training sequences
with that $v\_call$. Take the most common length, then compute the
per-position consensus (mode) amino acid. This gives a consensus
germline V-region $G_v$ per V gene, computed \emph{from the training
data itself} without external germline databases.
\item \textbf{Do the same for J genes}, using the FR4 segment after CDR3.
\item \textbf{Label each sequence.} For training sequence $x$ with
V-call $v$, J-call $j$, and CDR3 span $[a, b)$:
\begin{equation}
g_i = \begin{cases}
\Ind[x_i \neq G_{v,i}] & i < a \text{ (V-region)}, \\
0.5 & a \le i < b \text{ (CDR3 junction)}, \\
\Ind[x_i \neq G_{j, i - b}] & i \ge b \text{ (FR4)}.
\end{cases}
\end{equation}
\end{enumerate}

\subsection{Why CDR3 Gets $g_i = 0.5$}
The CDR3 junction is formed by V(D)J recombination with random nucleotide
insertions and is not derived from any single germline template. Its
residues are \emph{neither} germline-matched nor SHM-mutated in a
well-defined sense --- the very concept of ``germline'' is undefined for
junction positions. Assigning $g_i = 1$ (treating junction as fully
mutated) would over-concentrate mass on CDR3 and duplicate what CDR
masking already does. Assigning $g_i = 0$ would ignore CDR3 entirely.
The $0.5$ value puts CDR3 halfway, giving it moderate weight without
overlap with other strategies.

\subsection{Formal Definition}
With weights $(w_{\text{mut}}, w_{\text{germ}}) = (6, 1)$:
\begin{equation}
w^{(\text{germ})}_i = w_{\text{germ}} + (w_{\text{mut}} - w_{\text{germ}}) \cdot g_i
= 1 + 5 g_i.
\end{equation}
For a germline position ($g_i = 0$), $w_i = 1$. For a mutated
position ($g_i = 1$), $w_i = 6$. For CDR3 junction ($g_i = 0.5$),
$w_i = 3.5$. Masking probabilities via
Eq.~(\ref{eq:weighted-bernoulli}).

\subsection{Quantitative Consequence}
Assume a VH where $85\%$ of positions are germline-matched
($g_i = 0$, $w_i = 1$), $15\%$ are mutated ($g_i = 1$, $w_i = 6$),
and ignore CDR3 for simplicity. Then
\begin{equation}
\bar w = 0.85 \cdot 1 + 0.15 \cdot 6 = 0.85 + 0.90 = 1.75.
\end{equation}
The expected fraction of mask budget landing on mutated positions is
\begin{equation}
\frac{0.15 \cdot 6}{1.75} \approx 51\%,
\end{equation}
a $3.4\times$ boost over the $15\%$ rate a uniform masker would
produce. The expected fraction on germline positions is
$\frac{0.85 \cdot 1}{1.75} \approx 49\%$, still non-zero so the model
retains germline modeling capacity.

\subsection{Algorithm}
\begin{algorithm}[H]
\caption{Germline Masking}
\begin{algorithmic}[1]
\Require \texttt{germline\_labels} $g \in [0,1]^L$
\If{$g$ is \textsc{None}} \Return \Call{UniformMasking}{$x, s$} \EndIf
\State $w_i \gets w_{\text{germ}} + (w_{\text{mut}} - w_{\text{germ}}) g_i$
\State $w_i \gets 0$ for $s_i = 1$
\State $\bar w \gets \text{mean}\{w_i : s_i = 0\}$
\State $P_i \gets \min(\pmask w_i / \bar w, 1)$ if $\bar w > 0$ else $\pmask$
\State $P_i \gets 0$ for $s_i = 1$
\State $M \sim \Bern(P)$; \Return $M$
\end{algorithmic}
\end{algorithm}

\subsection{Design Justification}
\begin{itemize}
\item \textbf{Why build consensus from data instead of downloading IMGT
germline sequences?} (i) IMGT germlines are at the nucleotide level
and would require in-frame translation, handling of insertions/
deletions, and species-specific reference picks. (ii) A consensus
computed from the \emph{actual training data} is automatically
length-compatible and matches the tokenization. (iii) For any V gene
with many training sequences, the consensus $\to$ germline is
statistically indistinguishable from the true germline because SHM
is rare at any given position.
\item \textbf{Why length-group by the most common length?} V genes
occasionally have insertions/deletions within CDR1/2 that produce
multiple alignment lengths. Grouping by modal length gives a clean
set of aligned sequences for consensus; minority lengths fall back
to CDR-anchored segmented comparison (see \texttt{\_segment\_with\_cdrs}).
\item \textbf{Why min-seqs threshold of 20?} A consensus from $<20$
sequences is noisy; a single erroneous OAS record could flip a
consensus position. The threshold trades coverage (rare V genes) for
label quality. Sequences with rare V genes get $g_i = 0.5$ everywhere
(uncertain) and fall through to uniform-like behavior.
\item \textbf{Why $(w_{\text{mut}}, w_{\text{germ}}) = (6, 1)$ and not
extreme ratios like 100:1?} At $100{:}1$, with 15\% mutation rate,
virtually \emph{all} mask budget lands on the $\sim 18$ mutated
positions of a $\sim 120$-residue sequence, reducing to masking
``all mutations.'' This collapses the learning signal to a single
mode and destroys the budget-conservation property (all clamping).
$6{:}1$ produces a noticeable but controlled shift.
\item \textbf{Why the shared linear-interpolation form with
\texttt{interface.py}?} Both strategies have a single per-residue
label in $[0,1]$ and a weighted Bernoulli. The shared formula lets
hyperparameters (\texttt{mutated\_weight} vs \texttt{paratope\_weight})
be compared on the same footing.
\item \textbf{Why log the fallback rate?} Missing \texttt{germline\_labels}
is the most common silent failure (the sidecar path is optional).
Logging the fallback rate at \texttt{\_log\_interval = 1000} catches
this early.
\end{itemize}

% ============================================================================
\section{Strategy 7: Multispecific / Paired Chain Masking}
% ============================================================================

\subsection{Motivation}
Paired VH+VL antibodies (and multispecific constructs like bispecific
T-cell engagers) introduce problems that single-chain maskers cannot
address:
\begin{enumerate}
\item The VH--VL pairing carries information about binding --- which VL
can pair with which VH is not arbitrary.
\item Shared-light-chain bispecifics require conditional infilling:
given a fixed light chain, predict the heavy chain.
\item The VH--VL interface residues (Chothia/Martin canonical positions)
are a narrow set of cross-chain packing residues that determine
pairing geometry.
\end{enumerate}
Three policies are defined, one for each motivation, and sampled
stochastically at every training step.

\subsection{Required Metadata}
Every multispecific input must come with
\begin{itemize}[nosep]
\item $\mu_i \in \{0, 1, 2\}$: \texttt{module\_ids}, identifying which
module (global/mod1/mod2) token $i$ belongs to.
\item $\chi_i \in \{0, 1, 2\}$: \texttt{chain\_type\_ids}, where 0 is
special, 1 is heavy, 2 is light.
\end{itemize}
Optional: $\pi_i$ (\texttt{paratope\_labels}), $\iota_i$
(\texttt{interface\_labels}).

\subsection{Policy Mixture}
At each call, sample $k \sim \Cat(\rho_A, \rho_B, \rho_C)$ with
default $(\tfrac13, \tfrac13, \tfrac13)$, then apply policy $k$:

\subsubsection{Policy A: Module-Isolated Paratope Masking}
\paragraph{Idea.} Pick one module at a time; mask mostly its paratope;
leak a small fraction of the budget to the other module to prevent
the other-module representation from collapsing.

\paragraph{Formal definition.} Let $\mathcal{K}$ be the set of
non-empty modules. Sample a target $k^{\star} \sim \text{Uniform}(\mathcal{K})$.
Set
\begin{equation}
w^{(A)}_i =
\begin{cases}
w_{\bar{\text{p}}} + (w_{\text{p}} - w_{\bar{\text{p}}}) \pi_i & \mu_i = k^{\star},\, \pi \text{ present}, \\
w_{\text{p}} & \mu_i = k^{\star},\, \pi \text{ absent}, \\
0.1 \cdot w_{\bar{\text{p}}} & \mu_i \neq k^{\star},\, s_i = 0, \\
0 & s_i = 1.
\end{cases}
\end{equation}
Then feed $w^{(A)}$ through Eq.~(\ref{eq:weighted-bernoulli}).

\paragraph{The $0.1$ leak factor.} Without the leak ($0 \cdot w_{\bar{\text{p}}}$),
the non-target module would \emph{never} be masked under Policy A, and
its representation would gradually collapse toward an identity function
because there is no MLM signal on those tokens. With a $10\%$ leak, the
non-target module still receives $\approx 10\% \cdot \tfrac{n_{\text{other}}}{n_{\text{total}}}$
of the mask budget --- enough to prevent collapse but small enough to
preserve the module-isolated training signal.

\subsubsection{Policy B: Shared-Chain (Light-Chain) Boost}
\paragraph{Idea.} In bispecific antibodies, the light chain is commonly
shared between the two heavy chains; training the model to predict the
light chain conditioned on the heavy chain is the direct analogue of the
downstream ``design a light chain for this VH'' task. Even in
monospecific paired data, this teaches conditional VL$\mid$VH
inference.

\paragraph{Formal definition.}
\begin{equation}
w^{(B)}_i =
\begin{cases}
w_{\text{shared}} = 3 & \text{if } \chi_i = 2 \text{ (light)}, \\
1 & \text{if } \chi_i = 1 \text{ (heavy)}, \\
0 & s_i = 1.
\end{cases}
\end{equation}
With the budget normalization of Eq.~(\ref{eq:weighted-bernoulli}), the
expected mask budget on the light chain is
\begin{equation}
\frac{w_{\text{shared}} f_{L}}{w_{\text{shared}} f_{L} + f_{H}},
\end{equation}
where $f_L, f_H$ are the light/heavy length fractions.
For typical $f_L \approx f_H \approx 0.5$, this gives $\tfrac{3}{4}
= 75\%$ of the budget on light chain.

\subsubsection{Policy C: VH--VL Interface}
\paragraph{Idea.} Canonical VH--VL interface residues (Chothia \& Lesk
1985, Abhinandan \& Martin 2010) are a small set of positions
(\texttt{VH\_INTERFACE\_IMGT = \{37, 39, 45, 47, 91, 93, 100, 103\}},
\texttt{VL\_INTERFACE\_IMGT = \{36, 38, 44, 46, 87, 89, 98, 100\}}) that
form the direct cross-chain packing surface. Masking these residues
forces the model to use the opposite chain as context, which is the
cleanest possible signal for learning pairing.

\paragraph{Formal definition.} Given precomputed interface labels
$\iota_i \in [0, 1]$:
\begin{equation}
w^{(C)}_i = w_{\text{non-int}} + (w_{\text{int}} - w_{\text{non-int}}) \cdot \iota_i
= 1 + 5 \iota_i.
\end{equation}
$\iota_i$ is 1.0 at canonical interface positions, $0.3$ at
sequence-adjacent positions (one residue on either side), and 0
elsewhere.

\paragraph{Why the $0.3$ neighbor label?} The canonical interface
positions are a sparse set of $\sim 8$ residues per chain. Masking
\emph{only} these positions is statistically thin --- many training
samples would have only 1--2 interface residues. Including adjacent
residues at a lower weight (a) smooths the signal, (b) captures the
physical reality that packing is a continuous contact surface not an
isolated set of atoms, and (c) prevents Policy C from starving the
mask budget when $\sum \iota_i$ is very small.

\subsection{Multispecific Algorithm (Top Level)}
\begin{algorithm}[H]
\caption{Multispecific Masking}
\begin{algorithmic}[1]
\Require $\mu$, $\chi$; optionally $\pi, \iota$
\If{$\mu$ or $\chi$ missing} \Return \Call{UniformFallback}{} \EndIf
\State $k \sim \Cat(\rho_A, \rho_B, \rho_C)$
\If{$k = A$} \Return \Call{PolicyA}{$\mu, \chi, \pi$} \EndIf
\If{$k = B$} \Return \Call{PolicyB}{$\chi$} \EndIf
\If{$k = C$}
  \If{$\iota$ missing} \Return \Call{UniformFallback}{} \EndIf
  \State \Return \Call{PolicyC}{$\iota$}
\EndIf
\end{algorithmic}
\end{algorithm}

\subsection{Design Justification}
\begin{itemize}
\item \textbf{Why three separate policies and not a single weighted
union?} A weighted union would average out the three distinct
training signals: paratope, shared-chain, and interface. With
per-step policy sampling, each training step produces a mask that is
strongly committed to one of the three signals, so gradients remain
``peaky'' and distinguishable. Empirically (see AbLang-2,
ImmunoMatch), stochastic mixtures of distinct training regimes
outperform single-objective merges.
\item \textbf{Why equal default weights $(\tfrac13, \tfrac13, \tfrac13)$?}
No a priori reason to prefer one policy without downstream evidence.
Equal weights are the safest default for a first training run;
\texttt{policy\_weights} is exposed in the config for per-experiment
tuning.
\item \textbf{Why do Policy A's non-target modules get
$0.1 \cdot w_{\bar{\text{p}}}$ rather than $0.1 \cdot w_{\text{p}}$?}
The non-target module is not supposed to be the focus of the signal;
multiplying by the \emph{non-paratope} weight (not the paratope
weight) ensures that the leak is genuinely small, not a disguised
boost.
\item \textbf{Why Policy B boosts only the light chain and not the
heavy chain?} In bispecific formats (common-light-chain or CrossMab),
the light chain is the \emph{shared} component across different
antigens. Training the model to recover the light chain given the
heavy chain is exactly the design task. The reverse (mask heavy,
given light) is already well-covered by Policy A (paratope) and
Policies of single-chain strategies.
\item \textbf{Why fall back to uniform when interface labels are missing
under Policy C?} Policy C is meaningless without $\iota$; falling
through to uniform on that \emph{step only} preserves training
throughput.
\item \textbf{Why the hard interface position set from
Chothia/Martin?} These are the classical crystallographically-derived
interface positions --- they are canonical, well-documented, and
species-agnostic for human-like VH/VL. Alternative sources
(e.g.\ predicted interface from a docking model) would introduce
additional noise and dependencies.
\end{itemize}

% ============================================================================
\section{Strategy 8: Hybrid Masking (Meta-Strategy)}
% ============================================================================

\subsection{Motivation}
Each of the single-focus strategies (CDR, span, structure, interface,
germline) captures one aspect of antibody biology in isolation. But
the real pretraining objective is a model that understands all of
them simultaneously. Training on one strategy and then another
(naive curriculum) causes catastrophic interference; training on their
arithmetic union dilutes the signal.

The hybrid meta-strategy instead samples \emph{one sub-strategy per
training sample} from a configurable mixture distribution. Each
sample's mask is committed to a single strategy, so gradients per
sample are sharp, but over a batch the update is a mixture of all
strategies' gradients. This is the classical ``mixture of experts at
the loss level'' pattern.

\subsection{Formal Definition}
Let $\mathcal{S} = (S_1, S_2, \ldots, S_J)$ be the ordered tuple of
sub-strategies, and let $\pi^{(t)} \in \Delta^{J-1}$ be the mixture
distribution at training step $t$ (living on the probability simplex).
For each sample, the hybrid strategy:
\begin{enumerate}
\item Computes the \emph{available} sub-strategies for this sample:
$\mathcal{A} = \{j : S_j\text{ has all required metadata in this sample}\}$.
\item Renormalizes the mixture restricted to $\mathcal{A}$:
\begin{equation}
\tilde{\pi}^{(t)}_j = \frac{\pi^{(t)}_j \Ind[j \in \mathcal{A}]}{\sum_{j' \in \mathcal{A}} \pi^{(t)}_{j'}}.
\end{equation}
\item Samples $k \sim \Cat(\tilde{\pi}^{(t)})$ and delegates to
$S_k$.\texttt{select\_mask\_positions()}.
\end{enumerate}

\subsection{Metadata-Availability Filtering}
Each sub-strategy declares its required metadata keys in a static
table:
\begin{center}
\begin{tabular}{ll}
\toprule
Strategy & Required keys \\
\midrule
\texttt{uniform} & (none) \\
\texttt{span} & (none) \\
\texttt{cdr} & \texttt{cdr\_mask} \\
\texttt{structure} & \texttt{knn\_indices} \emph{or} \texttt{coords\_ca} \\
\texttt{interface} & \texttt{paratope\_labels} \\
\texttt{germline} & \texttt{germline\_labels} \\
\texttt{multispecific} & \texttt{module\_ids}, \texttt{chain\_type\_ids} \\
\bottomrule
\end{tabular}
\end{center}
The filter is applied per sample (not per batch) because sparsity of
annotations is heterogeneous across samples.

\paragraph{Why per-sample, not per-batch?} In a batch of 32, typically
$\sim 28$ will have all annotations and $\sim 4$ will be missing one.
Per-batch filtering would force the entire batch to drop a strategy if
any single member lacks it. Per-sample filtering fully uses every
sample's metadata.

\paragraph{Why restrict to $\mathcal{A}$ instead of falling back to
uniform?} Falling back to uniform on a sample where \emph{some} rich
strategies are unavailable would throw away the remaining information
(e.g.\ CDR labels are still present but germline labels are not).
Renormalizing lets the hybrid use whatever signal is available.

\subsection{Curriculum Scheduling}
The hybrid supports a piecewise-linear curriculum over $\pi^{(t)}$.
The user specifies breakpoints
$\{(t_1, \pi_1), (t_2, \pi_2), \ldots, (t_B, \pi_B)\}$ with
$t_1 < t_2 < \ldots < t_B$, and $\pi^{(t)}$ is defined by:
\begin{equation}
\pi^{(t)} = \begin{cases}
\pi_1 & t \le t_1, \\
(1 - \alpha) \pi_b + \alpha \pi_{b+1} & t_b \le t < t_{b+1},\ \alpha = \frac{t - t_b}{t_{b+1} - t_b}, \\
\pi_B & t \ge t_B.
\end{cases}
\end{equation}
After interpolation, $\pi^{(t)}$ is re-normalized to sum to $1$.

\subsubsection{Example Curriculum (from
\texttt{hybrid\_curriculum\_medium.yaml})}
\begin{center}
\small
\begin{tabular}{@{}r r r r r r r l@{}}
\toprule
Step & uniform & cdr & span & struct. & interf. & germ. & Stage \\
\midrule
0 & 0.30 & 0.15 & 0.30 & 0.10 & 0.10 & 0.05 & 1 \\
6{,}250 & 0.10 & 0.30 & 0.10 & 0.15 & 0.25 & 0.10 & 2 \\
18{,}750 & 0.10 & 0.20 & 0.10 & 0.15 & 0.20 & 0.25 & 3 \\
40{,}000+ & 0.10 & 0.20 & 0.10 & 0.15 & 0.20 & 0.25 & 4 (frozen) \\
\bottomrule
\end{tabular}
\end{center}
\vspace{0.5em}

\noindent \textbf{Stage rationale.} Stage 1 (general syntax) is
dominated by easy strategies (uniform + span) to let the model learn
basic protein language. Stage 2 (binding-site specialization) shifts
mass toward CDR and interface. Stage 3 (mutation realism) increases
germline weight as the model is now sophisticated enough to benefit
from the SHM signal. Stage 4 is frozen for the remaining $\sim 68\%$
of training at the balanced mixture.

\subsection{The Shared-Memory Curriculum Bug (and Fix)}
\paragraph{Problem.} When \texttt{dataloader\_num\_workers > 0},
PyTorch forks the data loader; each worker gets its own copy of the
\texttt{HybridMasking} object at fork time. A main-process callback
that updates \texttt{self.\_current\_weights = new\_weights} would
rebind the main-process attribute but leave worker copies unchanged
--- the curriculum would silently be a no-op.

\paragraph{Fix.} On initialization,
\texttt{self.\_current\_weights = self.\_base\_weights.clone().share\_memory\_()}
places the weight tensor in shared memory. The callback then updates
it via \emph{in-place} \texttt{copy\_()}, which writes to the underlying
shared buffer. All worker processes see the update because they hold
references to the same shared-memory tensor.

\paragraph{Why this is subtle.} The bug is silent: training runs,
logs succeed, and the curriculum \emph{appears} to work (the main
process logs show the updated weights) but the actual masks produced
in workers are generated from the initial weights forever. The fix is
one line but it is the difference between a working curriculum and a
silent research bug.

\subsection{Algorithm (Top Level)}
\begin{algorithm}[H]
\caption{Hybrid Masking}
\begin{algorithmic}[1]
\State \textbf{On init:}
\State \hspace{1em} $\pi^{(0)} \gets$ normalize(\texttt{policy\_weights})
\State \hspace{1em} $\pi^{(t)}$-buffer $\gets$ \texttt{share\_memory\_()}
\State \hspace{1em} Instantiate $S_1, \ldots, S_J$ via registry with sub-strategy params
\State \textbf{On step begin (callback):}
\State \hspace{1em} Interpolate $\pi^{(t)}$ from the curriculum; \texttt{copy\_()} into buffer
\State \textbf{On apply (per sample):}
\State \hspace{1em} $\mathcal{A} \gets \{j : S_j \text{ available for this sample}\}$
\If{$\mathcal{A} = \emptyset$}
  \State \Return \Call{UniformFallback}{}
\EndIf
\State $\tilde{\pi} \gets$ renormalize $\pi^{(t)}$ to $\mathcal{A}$
\State $k \sim \Cat(\tilde{\pi})$
\State \Return $S_k.\text{select\_mask\_positions}(x, s, \text{metadata})$
\end{algorithmic}
\end{algorithm}

\subsection{Design Justification}
\begin{itemize}
\item \textbf{Why per-sample policy sampling instead of per-batch?}
Per-batch would mean all 32 samples in a batch use the same masking
strategy, coupling gradient direction at the batch level. Per-sample
decorrelates strategy selection from batch identity, giving the
optimizer a cleaner Monte-Carlo estimate of the mixture objective.
\item \textbf{Why a mixture objective rather than simple averaging of
per-strategy losses?} Averaging losses requires computing the forward
pass once per strategy per sample --- cost multiplies by $J$. Mixture
sampling is an unbiased estimator of the same expected loss
\begin{equation}
\E_{k \sim \pi}[\mathcal{L}_k] = \sum_j \pi_j \mathcal{L}_j
\end{equation}
at $1\times$ cost per sample. The variance is higher, but at batch
size $32 \times$ gradient-accum $2$ $=$ $64$ samples per update,
the Monte-Carlo average is well-behaved.
\item \textbf{Why piecewise linear curriculum instead of exponential or
step function?} Step functions introduce discontinuities in the
training loss that can destabilize optimization. Linear interpolation
is smooth, has only two hyperparameters per segment (start, end), and
is easy to visualize. Exponential or cosine schedules were considered
but add parameters with no clear advantage.
\item \textbf{Why freeze the final weights at step 40k instead of
continuing to interpolate?} At step 40k the model has settled into a
mixture regime; continuing to shift weights would introduce
distributional drift at the end of training when we want to stabilize.
\item \textbf{Why is uniform always in the mixture?} Because
$\mathcal{A}$ is defined per sample, a sample with no metadata would
otherwise fall through to the hard-coded uniform fallback (at the
cost of per-sample logging). Including uniform as an explicit
sub-strategy with small weight ($\sim 10\%$) guarantees the fallback
path is exercised and benchmarked against biased strategies.
\item \textbf{Why renormalize $\pi$ over available sub-strategies
instead of using exactly $\pi_j$?} Without renormalization, samples
missing metadata for some strategy would have ``dead mass'' (the
total probability over $\mathcal{A}$ would be $< 1$), and the
Categorical sampler would implicitly down-weight those samples.
Renormalization ensures every sample makes full use of the mask
budget.
\item \textbf{Why track per-strategy counts
(\texttt{\_strategy\_counts})?} For diagnostics: if one strategy is
being over-selected because another's metadata is systematically
missing, the counts diverge from $\pi$ and the discrepancy shows up
in logs.
\end{itemize}

% ============================================================================
\section{Cross-Cutting Algorithmic Choices}
% ============================================================================

\subsection{Registry Pattern}
Strategies are registered via a \texttt{@register\_strategy("name")}
decorator into a module-level \texttt{\_STRATEGY\_REGISTRY} dict. The
training code resolves strategies by name string (from YAML) via
\texttt{get\_strategy()}. This decoupling means (a) adding a new
strategy requires zero changes to training code, (b) hybrid can
instantiate its sub-strategies by name without circular imports, and
(c) ablations are driven by config changes alone.

\subsection{Why All Strategies Share a Two-Stage Design}
The two-stage split (Stage 1: select positions; Stage 2: corrupt)
means every strategy only has to implement
\texttt{select\_mask\_positions()}. The BERT-style 80/10/10 corruption
is shared and consistent across strategies, so any ablation across
strategies is clean: the \emph{only} difference between runs is
\emph{which positions are masked}, not \emph{how they are corrupted}.

\subsection{Per-Sample Metadata Flow (Collator)}
The \texttt{MLMDataCollator} is strategy-agnostic. It:
\begin{enumerate}
\item Pads all sequences in a batch to $\max L$ (rounded up to a
multiple of 8 for Tensor Core alignment).
\item Builds per-sample metadata tensors by padding the sample's own
metadata to match (never re-indexing across the batch, which was a
historical bug).
\item Zero-fills missing metadata keys with reference shape \emph{and
dtype} (critical: float-typed metadata like
\texttt{paratope\_labels} must be zero-filled as float, not long).
\item Calls \texttt{strategy.apply(input\_ids, special\_mask,
metadata)} on each sample in turn.
\end{enumerate}

\paragraph{Why per-sample metadata instead of per-batch?}
Strategies see only the metadata relevant to this sample. A single
sample missing its \texttt{germline\_labels} does not poison the
entire batch's germline processing. This also means a strategy can
fall back cleanly on a per-sample basis.

\subsection{Determinism and Reproducibility}
All random draws use \texttt{torch.rand}, \texttt{torch.randint}, and
\texttt{torch.bernoulli}, so they are tied to the PyTorch RNG state.
A single seed at training-start fixes every strategy's random
behavior, making ablations reproducible.

% ============================================================================
\section{Summary Table}
% ============================================================================

\begin{center}
\renewcommand{\arraystretch}{1.3}
\small
\newcolumntype{L}{>{\raggedright\arraybackslash}X}
\begin{tabularx}{\textwidth}{@{}l L L L@{}}
\toprule
Strategy & Signal source & Position-selection rule & Key hyperparameters \\
\midrule
\texttt{uniform} & none & $p_i = \pmask$ (non-special) & $\pmask = 0.15$ \\
\texttt{cdr} & CDR region labels & Weighted Bernoulli, $(1,3,3,6)$ & $w_{\text{cdr1,2,3}}$ \\
\texttt{span} & none & Seed + Geom($p_g$) spans & $p_g = 0.2,\ \ell_{\max} = 10$ \\
\texttt{structure} & ESM2 contacts (kNN) & Seed-and-grow on kNN graph & $K = 32$ \\
\texttt{interface} & Soft paratope labels (TDC) & Weighted Bernoulli, $1 + 5\pi_i$ & $w_{\text{p}} = 6$ \\
\texttt{germline} & Consensus germline labels & Weighted Bernoulli, $1 + 5 g_i$ & $w_{\text{mut}} = 6$ \\
\texttt{multispecific} & Module, chain, paratope, interface & 3-policy mixture (A/B/C) & $w_{\text{shared}} = 3$ \\
\texttt{hybrid} & All of the above & Per-sample Categorical mixture with curriculum & $\pi^{(t)}$ schedule \\
\bottomrule
\end{tabularx}
\end{center}

% ============================================================================
\section{Likely Questions (and Answers)}
% ============================================================================

\paragraph{Q: Why a fixed $\pmask = 0.15$ across all strategies, rather
than tuning per-strategy?}
Keeping $\pmask$ fixed is the only way to isolate the
position-selection axis in ablations. Two strategies with different
budgets produce different total gradient norms, confounding the
comparison. Tuning $\pmask$ is an orthogonal study.

\paragraph{Q: How do you know weighted Bernoulli conserves the expected
budget?}
See Eq.~(\ref{eq:budget-conservation}). The normalization by $\bar w$
is exactly the factor that makes $\E[\sum M_i]/n = \pmask$ before
clamping. Empirically we checked the realized mask rates on held-out
batches and they match $\pmask$ within $\pm 0.5\%$.

\paragraph{Q: What if a sequence has 0 CDR residues (annotation failed)?}
Then $\bar w = w_{\text{fr}} = 1$ and Eq.~(\ref{eq:weighted-bernoulli})
reduces to $p_i = \pmask \cdot 1/1 = \pmask$, i.e.\ exact uniform
masking. The strategy degrades gracefully to the baseline without
any special-casing, which is a property of the weighted-Bernoulli
formulation.

\paragraph{Q: Why not use a single ``universal'' weight vector that
combines CDR, germline, paratope, and structure?}
Multiplicative combination (e.g.\ $w_i = w^{(\text{cdr})}_i \cdot w^{(\text{germ})}_i$)
produces an extreme distribution where the top few positions get the
bulk of the probability. Additive combination
($w_i = w^{(\text{cdr})}_i + w^{(\text{germ})}_i$) produces a blurred
signal that is dominated by whichever strategy has the largest
magnitude. The hybrid sampler sidesteps both failure modes by
committing to one signal per sample.

\paragraph{Q: Why not learn the mixture weights $\pi$?}
A learned $\pi$ couples the masker to the model and creates a
feedback loop: if the model is bad at paratope prediction, gradient
descent would down-weight paratope masking, further starving paratope
learning. Fixed weights (or a hand-crafted curriculum) maintain a
stable learning target.

\paragraph{Q: How do you handle sequences shorter than a span?}
Span masking stops at the sequence boundary ($t + o \ge L$) and
special tokens. If $T$ cannot be filled (e.g.\ a very short sequence
with no room), the \texttt{max\_attempts} cap breaks out after $10T$
tries. The resulting mask is below budget but still valid; the
gradient contribution is just proportionally smaller.

\paragraph{Q: What about bias introduced by the soft paratope labels
being derived from TDC data (which is small and structurally biased)?}
This is an acknowledged limitation. The empirical table is smoothed
(region-level fallback for low-count cells) and the labels are used
as \emph{weights}, not hard targets --- a biased weight still
concentrates mass without lying about ground truth. Hard paratope
labels from crystal structures would be cleaner but are unavailable
at the 500k scale.

\paragraph{Q: Why does the structure strategy not mask by sequence
distance too?}
Sequence-distance masking is exactly what span masking does. The
structure strategy is specifically the \emph{complement} ---
coupling sequence-distant residues that are close in 3D. Combining
them would collapse to something close to uniform.

\paragraph{Q: How does the hybrid curriculum interact with the
learning rate warmup?}
The LR warmup runs $0 \to t_{\text{warm}} = 6{,}250$ steps, which
coincides with Stage 1 $\to$ Stage 2 of the curriculum. So the
easiest mixture (mostly uniform + span) is used during warmup when
the model is still unstable, and the harder biased mixtures kick in
as the learning rate reaches its plateau. This is intentional: we
want the harder training signal only when the optimizer is ready for it.

\end{document}
